\chapter{Sprint 6: AI Copilot \& Estimation (RAG)}

Five sprints built the platform: identity, projects, the agile and task engines, the productivity suite, and
the communication and monitoring backbone. Since Sprint~1, the platform has been accumulating one thing that
this last sprint finally puts to work: \emph{project data}. Tasks, comments, epics, sprints, and milestones
together form a written record of how the team actually works: what was decided, what is blocking whom, how
long a piece of work really took. Sprint 6 turns that record into two assistive features without installing a
separate machine-learning stack: a copilot that answers natural-language questions grounded strictly in the
team's own content, and an estimator that bases a draft task's effort on the measured outcomes of similar
completed work. It is the project's main technical contribution and, unlike the earlier sprints, it comes
with quantitative evidence that it works. For this reason, this chapter treats it in more depth.

\section{Sprint goal}

\textbf{Sprint goal.} Give the platform two features that read its own data back to the user. The first is a
\textbf{copilot} that answers questions like ``why did we choose X?'' or ``what is blocking Y?'' grounded only in
retrieved project content, with clickable citations back to the exact task or comment, and that refuses
honestly when the corpus cannot support an answer. The second is an \textbf{estimation} feature that predicts a
draft task's effort from the real \texttt{actualHours} of the most similar completed tasks. Both are strictly
permission-scoped: a user can only ever retrieve from projects they may access. One design principle runs
through the whole subsystem, \emph{honesty over coverage}: refuse rather than guess, show an uncertainty band
instead of a false point estimate, and cite the source instead of just stating it.

\textbf{Why RAG, and why not a bare LLM.} The naive design --- sending the question straight to a large language
model --- fails for exactly the reason this platform exists: the model knows nothing about \emph{this} company's
projects, and asked about them, it will invent a plausible answer. Retrieval-Augmented
Generation~\cite{lewis2020rag} fixes that by reversing the flow. We first \emph{retrieve} the few real project
chunks most relevant to the question, then give only those to the model as its sources, with an instruction to
answer from them and cite them by number. The model becomes a summarizer of retrieved evidence rather than an
oracle. An answer can therefore always be traced back to a source the user can open. Most importantly, when
nothing relevant is retrieved, the system can \emph{decline} instead of inventing. We enforce that decline with
a confidence gate: if the best retrieved chunk is not similar enough to the question, the copilot refuses
without ever calling the generation model.

The same logic drives estimation. Rather than ask an LLM to guess ``how many hours will this
take?'' --- a question it has no real basis to answer --- we treat estimation as \textbf{reference-class
forecasting}~\cite{kahneman1977reference}: find the completed tasks most similar to the draft, and base the estimate on how long \emph{they}
actually took. The suggestion is backed by evidence and comes with the neighbour tasks attached, so the user
sees where the number comes from and can judge it. In both features, the point is not that AI is
present. The value is that its output stays grounded in real data, scoped to a user's permissions, and honest
about what it does not know.

\section{Sprint backlog}

The stories below are drawn from the product backlog (Section~2.2). At \textbf{42 story points} it is a large
sprint even though it touches only one module, because that module is deep: an asynchronous indexing pipeline,
a hybrid retrieval engine, a grounded-generation orchestrator, an estimator, and an evaluation harness.

\begin{longtable}{@{}l>{\raggedright\arraybackslash}p{0.50\textwidth}lll@{}}
\caption{Sprint 6 backlog.}\label{tab:s6-backlog-detail} \\
\tableheaderrow
\thcell{US-ID} & \thcell{Story} & \thcell{SP} & \thcell{Est.} & \thcell{Priority} \\
\midrule
\endfirsthead
\caption[]{Sprint 6 backlog (continued).} \\
\tableheaderrow
\thcell{US-ID} & \thcell{Story} & \thcell{SP} & \thcell{Est.} & \thcell{Priority} \\
\midrule
\endhead
\bottomrule
\endfoot
US-S6-01 & As a user, I want to ask the copilot a natural-language question about project content, so that I get answers without manual digging. & 5 & 2.5d & Must \\
US-S6-02 & As a user, I want the answer streamed with clickable citations to the source items, so that I can trust and verify it. & 8 & 4d & Must \\
US-S6-03 & As a user, I want the copilot to refuse honestly when the corpus cannot support an answer, so that I am never misled by a hallucination. & 5 & 2d & Must \\
US-S6-04 & As a member, I want the system to estimate a draft task's effort from similar completed tasks, so that planning is data-driven. & 8 & 3.5d & Should \\
US-S6-05 & As the platform, I want project content indexed automatically via a write-path outbox, so that the copilot stays current without slowing saves. & 8 & 4d & Must \\
US-S6-06 & As an admin, I want to trigger a reindex and view copilot telemetry, so that I can operate the AI subsystem. & 3 & 1.5d & Should \\
US-S6-07 & As the platform, I want retrieval scoped by permissions in SQL, so that no user can surface content from a project they cannot access. & 5 & 2.5d & Must \\
\midrule
 & \textbf{Sprint 6 subtotal} & \textbf{42} & \textbf{4w} & \\
\end{longtable}

\section{Requirements analysis}

The sprint delivers a single AI module with two user-facing features --- a copilot and a task
estimator --- sharing the same indexing infrastructure and retrieval engine.

Figure~\ref{fig:4-1} shows the features, actors, and use cases. A plain \textbf{User} (holding \texttt{TASKS.TASK\_READ\_MANY}) asks the copilot; a \textbf{task
author} (holding \texttt{TASKS.TASK\_CREATE}) triggers an estimate while drafting a task; an \textbf{admin / executive}
(holding \texttt{PROJECTS.PROJECT\_CREATE}) operates the subsystem by reindexing and reading telemetry. A \textbf{Scheduler}
actor represents the automated indexer --- no human triggers embedding; content flows into the index
automatically after a write. The two user-facing features both \texttt{<<include>>} a permission-scoped \emph{retrieve} use case, and
an honest refusal \texttt{<<extend>>}s the streamed, cited answer when retrieval cannot support one.

% Figure 4.1 — render diagrams-src/fig_4_1.puml -> figures/fig_4_1.(png|pdf)
\reportfigure{figures/fig_4_1}{Sprint 6 use-case diagram with RBAC-labelled actors.}{4-1}

\subsubsection{Module A --- Copilot}

\textbf{Functional description.} The copilot lets any team member ask a natural-language question
about their project's content --- tasks, comments, epics, milestones, and sprints --- and receive
a streamed answer grounded strictly in retrieved project data, with clickable citations back
to the exact source items. When the corpus cannot support an answer, the copilot refuses
honestly rather than generating a plausible but unfounded response. An admin can trigger a
full reindex and view copilot telemetry. The indexing pipeline runs asynchronously via an
outbox pattern: content changes are enqueued alongside their domain writes and processed by a
background sweeper, so a failure in the AI subsystem can never break a task save.

\textbf{Actors.} A \textbf{User} (\texttt{TASK\_READ\_MANY}) asks questions and receives grounded
answers. An \textbf{Admin / Executive} (\texttt{PROJECT\_CREATE}) triggers a reindex and views
telemetry. A \textbf{Scheduler} (system) runs the indexing sweeper and nightly reconcile
automatically.

\textbf{Business rules.}
\begin{itemize}
  \item Retrieval is \textbf{permission-scoped in SQL}: a user can only retrieve content from
    projects they may access; out-of-scope rows are never returned to rank, making
    cross-project leakage impossible by construction rather than merely checked.
  \item A \textbf{confidence gate} refuses without calling the generation model when the best
    retrieved chunk is insufficiently similar to the question (cosine $< 0.5$).
  \item Answers are streamed with numbered \textbf{citation markers} (\texttt{[n]}) linking back
    to the source entity; each citation resolves to a deep-link the user can open.
  \item The system instruction and retrieved (untrusted) sources are passed in
    \textbf{separate turns}, establishing a prompt-injection boundary.
  \item Five entity types are indexed: tasks, task comments, epics, milestones, and sprints.
  \item The outbox has a \texttt{@@unique(entityType, entityId)} constraint, \textbf{collapsing
    repeated edits} to a single pending row.
  \item A nightly reconcile re-queues any source newer than its embedding and enqueues deletes
    for orphans, so the index \textbf{self-heals} from any missed write.
\end{itemize}

\subsubsection{Module B --- Task Estimation}

\textbf{Functional description.} The estimation feature predicts a draft task's effort from the
measured outcomes of similar completed tasks, treating effort prediction as reference-class
forecasting~\cite{kahneman1977reference} rather than asking a language model to guess. It reuses the copilot's
scoped retrieval to find the nearest completed tasks by semantic similarity, and reports a
data-driven suggestion with the evidence tasks attached, so the user sees the basis for the
number and can judge it. No generation model is called --- the feature is pure retrieval.

\textbf{Actors.} A \textbf{Task Author} (\texttt{TASK\_CREATE}) triggers an estimate while
drafting or editing a task.

\textbf{Business rules.}
\begin{itemize}
  \item Only completed tasks with recorded actual hours (\texttt{actualHours > 0}) are considered
    as neighbours.
  \item If the target project yields fewer than \textbf{three neighbours}, the search widens to
    the same business unit.
  \item Two prediction modes: \textbf{size-aware} (when the draft and neighbours carry story
    points) computes draft points $\times$ a similarity-weighted hours-per-point percentile,
    reported as a median with a 10th/90th band; \textbf{size-agnostic} (fallback) computes a
    similarity-weighted median of \texttt{actualHours} with a 25th/75th IQR band.
  \item The suggestion is \textbf{never a bare number}: the neighbour tasks are always returned
    as evidence.
\end{itemize}

\section{Methodology note --- a CRISP-DM sub-process}

Scrum organizes the project as a whole, but the AI work inside this sprint is a data-mining task, and we
organized it using the vocabulary of \textbf{CRISP-DM}~\cite{shearer2000crisp} (Cross-Industry Standard Process for Data Mining) as a
sub-process inside Sprint 6. The mapping is direct, and it explains the shape of the subsystem:

\begin{itemize}
  \item \textbf{Business understanding.} The goal is assistive, not autonomous: help a user find and estimate, never
  decide for them. This is what makes \emph{honesty over coverage} the governing rule.
  \item \textbf{Data understanding.} The corpus is the platform's own content --- five entity types (tasks, task comments,
  epics, milestones, sprints) --- small, clean, permission-partitioned, and already relational. There is no
  external dataset to clean; the data understanding step is really \emph{scope} understanding (who may see what).
  \item \textbf{Data preparation.} Chunking (\textasciitilde2000-character windows with 300-character overlap, per entity type), sha256
  hashing to skip unchanged content, and embedding into 1536-dimensional vectors. This is the indexing
  pipeline of Section~4.5.
  \item \textbf{Modeling.} Not a trained model but a \emph{retrieval} model: hybrid vector + lexical search fused by Reciprocal
  Rank Fusion~\cite{cormack2009rrf}, an optional LLM reranker, and a k-NN estimator over completed-task outcomes.
  \item \textbf{Evaluation.} An offline harness with committed gold sets measuring Recall@k, MRR, nDCG, faithfulness,
  citation precision/recall, refusal correctness, and estimation error --- the subject of Section~4.8.
\end{itemize}

We did not build a \emph{deployment} step in the CRISP-DM sense of shipping a model artefact, because there is no
model to ship: the ``model'' is Google Gemini~\cite{gemini} behind an API, and our contribution is the retrieval,
grounding, and scoping around it. Naming these phases gives the work an explicit structure --- it makes clear
that the evaluation harness is the quality gate that plays the role of a trained model's validation curve.

\section{RAG architecture}

The subsystem follows the project's layering convention (controller \textrightarrow{} service \textrightarrow{} repository), with two
Prisma-raw repositories isolating the pgvector~\cite{pgvector} and tsvector SQL that Prisma cannot type. It is three
tables, five endpoints, nine single-responsibility services, and one cron worker carrying two schedules (the
per-minute sweep and the nightly reconcile). Figure~\ref{fig:4-2} presents the
whole pipeline in one view; the two sequence diagrams and the class diagram that follow detail the write path,
the read path, and the data model.

% Figure 4.2 — render diagrams-src/fig_4_2.puml -> figures/fig_4_2.(png|pdf)
\reportsideways{figures/fig_4_2}{RAG pipeline overview.}{4-2}

\textbf{The indexing pipeline (write path).} The central design choice is that \textbf{no Gemini call ever happens on a
user's request}. When a task is created, updated, or deleted, the task service calls a one-line
\texttt{enqueueUpsert} / \texttt{enqueueDelete} into an \textbf{outbox} table alongside its other side effects, and returns
immediately. The enqueue is deliberately fire-and-forget --- its failures are caught and logged, never
propagated --- so a problem in the AI subsystem can never break a task save. A separate
\texttt{IndexSweeperJob} cron runs every minute under a Postgres distributed lock (so only one instance sweeps, the
same pattern as the reminder scheduler), claims up to 25 due rows, and for each one re-reads the \emph{live} source
row, rebuilds its chunks, embeds only the chunks whose sha256 hash changed, and upserts them. Success marks
the row done; a transient failure schedules a retry with exponential backoff (30 s up to \textasciitilde15 min); when
the retry budget is exhausted, the row is marked failed. A nightly reconcile re-queues any source newer than its embedding and
enqueues deletes for orphaned embeddings, so the index self-heals from any missed write.
Two properties make this safe under repeated edits: the outbox has a \texttt{@@unique(entityType, entityId)}
constraint, so a task edited ten times before the next sweep collapses to a single pending row and is embedded
once; and re-reading the live row at sweep time means a stale queued edit always embeds the \emph{current} content.
Figure~\ref{fig:4-3} follows one task through this path.

% Figure 4.3 — render diagrams-src/fig_4_3.puml -> figures/fig_4_3.(png|pdf)
\reportsideways{figures/fig_4_3}{Asynchronous index-a-task sequence.}{4-3}

\textbf{The retrieval and answer pipeline (read path).} A copilot question runs through \texttt{RetrievalService}, which
first resolves the user's \texttt{allowedProjectIds} (an executive CEO sees all projects; a CTO sees the Tawer Dev
unit, a CMO Tawer Creative; everyone else sees only projects they are a member of), then embeds the question
as a retrieval query and runs two searches under the identical permission filter. The \textbf{vector arm}
(\texttt{searchVector}) does a cosine approximate-nearest-neighbour scan over the HNSW index; the \textbf{lexical arm}
(\texttt{searchLexical}) does a Postgres full-text \texttt{ts\_rank\_cd} scan over the GIN index. Their two ranked lists are
fused by \textbf{Reciprocal Rank Fusion}~\cite{cormack2009rrf} --- each candidate scores \texttt{$\Sigma$ 1/(60 + rank)} across the arms, which rewards
items ranked highly by either signal without needing the two scores to be on the same scale. An optional LLM
reranker can re-order the fused pool, and falls back safely to the fused order if the model is unavailable.

Then the \textbf{confidence gate}: if the best cosine score is below 0.5, retrieval reports \texttt{sufficient: false} and
the copilot refuses \emph{without ever calling the generation model} --- the cheapest possible way to maintain honesty.
Otherwise the service builds a numbered-sources prompt, passing the system instruction and the retrieved
(untrusted) sources in \textbf{separate turns} so injected instructions inside project text are treated as data,
and streams a grounded answer at low temperature (0.2). It parses the \texttt{[n]} citation markers the model emits
back to their source entities, batch-resolves their human-readable labels and deep-links, and writes a
telemetry receipt. Figure~\ref{fig:4-4} traces the streaming path.

% Figure 4.4 — render diagrams-src/fig_4_4.puml -> figures/fig_4_4.(png|pdf)
\reportsideways{figures/fig_4_4}{Copilot query sequence diagram.}{4-4}

\textbf{Estimation (read path, no generation).} Estimation reuses the same scoped retrieval but never calls the
generation model. It embeds the draft's \texttt{title} + \texttt{description}, runs a k-NN search restricted to \texttt{status = 'DONE' AND actualHours > 0} tasks (taking each task's single nearest chunk), and if the project yields fewer
than three neighbours widens the search to the same business unit. From the neighbours it predicts in two ways:
with draft story points present and pointed neighbours, it computes a \textbf{size-aware} estimate --- draft points \texttimes{}
a similarity-weighted hours-per-point percentile, reported as a median with a 10/90 band; otherwise it falls
back to a \textbf{size-agnostic} similarity-weighted median of \texttt{actualHours} with a 25/75 IQR band. Either way the
neighbours are returned as evidence, so the suggestion is never a bare number.

\textbf{The data model.} Three tables support the subsystem, modeled together in a dedicated Prisma schema file; the
vector column, the generated \texttt{tsvector} column, and the HNSW and GIN indexes are created in hand-written
migration SQL because Prisma cannot emit them. Figure~\ref{fig:4-5} presents the class diagram.
\texttt{DocumentEmbedding} holds one embedded chunk --- a \texttt{vector(1536)} column indexed by HNSW for cosine ANN, a
\emph{generated, stored} \texttt{tsvector} column indexed by GIN for full-text search (generated-and-stored means it
back-fills every existing row at migration time and self-maintains on write, with no write-path change), and a
\texttt{contentHash} so an unchanged chunk skips the Gemini~\cite{gemini} call. \texttt{IndexOutbox} is the async work queue with its
collapse-to-one \texttt{@@unique(entityType, entityId)} key and its backoff fields. \texttt{CopilotQueryLog} is one
telemetry receipt per copilot call. One structural note: \texttt{entityId} on the first two tables is a \textbf{soft
reference} --- a bare id column with no foreign key --- because it is polymorphic across the five entity types;
referential integrity is maintained by the write-path producers and the nightly reconcile rather than by the
database.

% Figure 4.5 — render diagrams-src/fig_4_5.puml -> figures/fig_4_5.(png|pdf)
\reportfigure{figures/fig_4_5}{AI module class diagram.}{4-5}

\section{Security considerations}

Security is a real strength of this module, and the reason is structural rather than procedural: the
protection comes from the query design, not from checks added afterwards.
\textbf{Retrieval permission scoping is enforced in SQL, before ranking.} \texttt{AiAccessService.allowedProjectIds}
resolves the caller's scope, and every search --- vector and lexical alike --- applies \texttt{projectId = ANY(\$allowedIds)} inside the query itself. There is no code path that can surface a chunk from a project the
user cannot access, because the out-of-scope rows are never returned to rank in the first place; the main
safety property is not just checked --- it is impossible by construction. An explicit out-of-scope
\texttt{projectId} returns 403. This is the property the evaluation independently confirmed with zero
cross-role leaks (Section~4.8).

Four further controls reinforce it:

\begin{itemize}
  \item \textbf{Endpoint RBAC} on every route via \texttt{HasPermissionGuard}: copilot requires \texttt{TASK\_READ\_MANY}, estimate
  \texttt{TASK\_CREATE}, and reindex/telemetry \texttt{PROJECT\_CREATE} (executive-only).
  \item \textbf{A prompt-injection boundary.} The system instruction and the untrusted retrieved sources are passed in
  separate turns, and the instruction explicitly tells the model to treat source text as data and never obey
  instructions embedded in it --- the same rule applies in the reranker. Temperature is held at 0.2 for faithful,
  low-variance answers.
  \item \textbf{DTO validation} bounds every input (question 3--1000 chars, story points 1--100, UUIDs).
  \item \textbf{No SQL-injection surface despite raw SQL.} Every raw query is a parameterized tagged template; the
  lexical arm parses user input with \texttt{websearch\_to\_tsquery}, which tolerates arbitrary text without throwing;
  the only interpolated value is the app-built numeric vector literal.
\end{itemize}

One limitation must be stated: access control is \textbf{project-level, not entity-level}. Any content in an
allowed project is retrievable, with no per-task or per-comment visibility check. That matches the tasks
module's own project-scoped authorization today, but if finer-grained task visibility were ever added,
retrieval would expose too much until it applied the same rule. This and the telemetry-privacy point
are carried into the sprint review (Section~4.10).

\section{Implementation}

On the frontend the copilot is a project-detail tab: a textarea and an ask button, the answer rendered
token-by-token behind a blinking caret as the SSE stream arrives, and then a row of citation chips --- or an
honest ``not found in this project's content'' note when retrieval was insufficient. The stream client is a
manual \texttt{fetch} + SSE-frame parser rather than the browser \texttt{EventSource}, specifically so it can send the
\texttt{Authorization} header, with one silent token-refresh retry on a 401. This is a deliberate engineering
choice: \texttt{EventSource} cannot set request headers, so a permission-scoped streaming endpoint behind JWT
auth would otherwise have to pass the token through the query string. Instead, the module reads the response
body as a stream and splits the SSE frames itself. The bearer token stays in the header, and a mid-stream 401
silently refreshes the token and replays the request, so the user never sees a dropped answer.
Each \textbf{citation chip} deep-links by
rewriting the URL's tab and task parameters: a task or comment citation opens the tasks tab and the task
sheet, while an epic, milestone, or sprint citation switches to the relevant tab, with the source snippet on
hover. Estimation appears inline in the create/edit task form as a single suggestion line --- ``$\approx$ Xh (low--high) $\cdot$
N pts --- based on TASK-a\ldots'' --- with a one-click apply, debounced and gated on the title reaching a few
characters, with appropriate empty and error states.

% Screenshot P-31
\reportscreenshot{figures/screenshot_P31}{Copilot panel: grounded answer with citation chips.}{P31}

% Screenshot P-32
\reportscreenshot{figures/screenshot_P32}{Citation chip deep-link opening a referenced task.}{P32}

% Screenshot P-33
\reportscreenshot{figures/screenshot_P33}{Task estimate suggestion in the create-task form.}{P33}

\section{Evaluation \& testing}

Unlike every other module, the AI subsystem comes with quantitative evidence. Rather than unit tests it
carries a substantial \textbf{offline evaluation harness} (\texttt{src/ai/eval/}, run via \texttt{npm run ai:eval:*}) with
committed gold sets and stable human-readable references, measuring three things: \textbf{retrieval} quality
(Recall@k, Precision@k, MRR, nDCG@k over 12 semantic and 10 keyword gold questions, comparing vector-only vs
hybrid vs hybrid+rerank, plus ablations); \textbf{answer} quality (an LLM faithfulness judge, citation
precision/recall against a \texttt{mustCite} set, and refusal correctness on answerable vs deliberately-unanswerable
questions); and \textbf{estimation} (leave-one-out MAE/RMSE and within-$\pm$25\% of the k-NN predictor against
project-mean and story-points baselines). This is measurement rather than regression testing, but it is the
de-facto quality gate that plays the role of the trained-model validation a CRISP-DM process would expect.
The three retrieval arms are themselves the ablation: \textbf{vector-only} (baseline),
\textbf{vector + lexical fused by RRF}~\cite{cormack2009rrf} (hybrid), and \textbf{hybrid followed by the LLM reranker} --- each flag-gated
(\texttt{AI\_RETRIEVAL\_HYBRID}, \texttt{AI\_RETRIEVAL\_RERANK}), so any arm can be switched off and re-measured in isolation.

The main result is on retrieval. On the \textbf{keyword / identifier} gold set (re-run 2026-07-07) --- bare
ticket-key lookups like \texttt{NDF-24}, the worst case for a dense embedder because an opaque token carries almost
no semantic signal --- adding the lexical arm changes the numbers completely:

\begin{longtable}{@{}lcccccc@{}}
\caption{Keyword/identifier gold-set retrieval metrics.}\label{tab:eval-keyword} \\
\tableheaderrow
\thcell{Config} & \thcell{MRR} & \thcell{R@1} & \thcell{nDCG@1} & \thcell{R@3} & \thcell{nDCG@3} & \thcell{R@5} \\
\midrule
\endfirsthead
\caption[]{Keyword/identifier gold-set retrieval metrics (continued).} \\
\tableheaderrow
\thcell{Config} & \thcell{MRR} & \thcell{R@1} & \thcell{nDCG@1} & \thcell{R@3} & \thcell{nDCG@3} & \thcell{R@5} \\
\midrule
\endhead
baseline (vector-only) & 0.570 & 0.300 & 0.300 & 0.900 & 0.639 & 1.000 \\
\textbf{hybrid} & \textbf{1.000} & \textbf{1.000} & \textbf{1.000} & \textbf{1.000} & \textbf{1.000} & \textbf{1.000} \\
hybrid + rerank & 1.000 & 1.000 & 1.000 & 1.000 & 1.000 & 1.000 \\
\bottomrule
\end{longtable}

That is \textbf{MRR 0.57 \textrightarrow{} 1.00 and Recall@1 0.30 \textrightarrow{} 1.00} at no extra LLM cost --- vector-only put the exact task
first only 30\% of the time, while the lexical arm matches the identifier every time and RRF fuses it to a
perfect ranking. On the \textbf{semantic} gold set the result is the opposite and equally informative: the
bi-encoder is already saturated on this small 73-chunk corpus (MRR 0.958, the first relevant hit at rank 1
almost every time), so \textbf{hybrid ties the saturated baseline (0.958, no regression); the LLM reranker matches
or slightly beats it (1.000 this run)}. Because the reranker is an LLM scoring pass, its number varies
slightly run-to-run, while the deterministic vector and hybrid arms do not --- so the accurate reading is that
rerank never regresses the semantic set and sometimes improves it. Hybrid is the right default: it is far
better in the keyword case and equal in the semantic one.

\begin{longtable}{@{}p{2.6cm}p{2.4cm}p{5.2cm}p{4.2cm}@{}}
\caption{Retrieval evaluation summary.}\label{tab:eval-summary} \\
\tableheaderrow
\thcell{Gold set} & \thcell{Metric} & \thcell{Result} & \thcell{Reading} \\
\midrule
\endfirsthead
\caption[]{Retrieval evaluation summary (continued).} \\
\tableheaderrow
\thcell{Gold set} & \thcell{Metric} & \thcell{Result} & \thcell{Reading} \\
\midrule
\endhead
Keyword (10 Q) & MRR / R@1 & 0.57 \textrightarrow{} \textbf{1.00} / 0.30 \textrightarrow{} \textbf{1.00} (hybrid vs vector) & hybrid closes the keyword gap \\
Semantic (12 Q) & MRR / R@1 / nDCG@1 & 0.958 / 0.583 / 0.917 (baseline = hybrid); rerank \textbf{1.000} / --- / \textbf{1.000} & embedder saturated; hybrid no regression, rerank slightly better \\
Cross-role leakage & out-of-scope hits & \textbf{0} in both arms & confirms the in-SQL scoping of Section~4.6 \\
\bottomrule
\end{longtable}

The cross-role leakage check ran the same query as a CEO (6 projects in scope) and as an intern (1 project)
and confirmed each actor's candidates came only from their allowed projects --- zero out-of-scope hits in either
arm, exactly as the shared \texttt{projectId = ANY(\$allowedIds)} filter predicts.

\textbf{What a retrieval miss looks like.} The vector-only failures on the keyword set are not random --- they have a
single, explainable cause. A bare identifier like \texttt{NDF-24} carries almost no semantic meaning, so the
dense embedder places it among many unrelated chunks and ranks the true task low (hence R@1 0.30). The lexical
arm has no such difficulty: \texttt{websearch\_to\_tsquery} matches the token exactly and returns the one chunk that
contains it. RRF then fuses the two ranked lists and promotes that exact-match hit to rank 1. The failure and
its fix are therefore structural, not the result of tuning --- the lexical arm exists exactly to cover the type
of queries the embedder cannot handle.

\textbf{Answer quality.} Answer quality was measured with the same hybrid retrieval, re-run on 2026-07-07. Over the
10-question QA gold set (7 answerable, 3 deliberately unanswerable), the copilot was fully grounded on every
judged answer (mean faithfulness 1.000), reached citation precision 0.786 and recall 1.000, and handled
refusal perfectly: all three unanswerable questions were correctly refused (correct-refusal rate 1.000) with a
zero false-refusal rate (actor CEO, 6 projects in scope).

\begin{longtable}{@{}p{9cm}p{3cm}@{}}
\caption{Answer-quality metrics.}\label{tab:eval-answers} \\
\tableheaderrow
\thcell{Metric} & \thcell{Result} \\
\midrule
\endfirsthead
\caption[]{Answer-quality metrics (continued).} \\
\tableheaderrow
\thcell{Metric} & \thcell{Result} \\
\midrule
\endhead
Mean faithfulness (grounded answers) & 1.000 \\
Citation precision / recall & 0.786 / 1.000 \\
Refusal accuracy (all 10) & 1.000 \\
Correct refusals (3 unanswerable) & 1.000 \\
False-refusal rate & 0.000 \\
\bottomrule
\end{longtable}

The faithfulness mean is over n = 6: the LLM judge returned no verdict on one otherwise correctly-cited
answer, which we report rather than silently drop.

\textbf{Estimation quality.} The estimator was evaluated by \textbf{leave-one-out} over the 43 completed tasks with
recorded hours: each task is held out in turn, predicted from its neighbours, and compared to its real
\texttt{actualHours}. Two baselines put the numbers in context: the \emph{project mean} (predict every task at its project's
average) and a \emph{story-points fit} (a least-squares hours-per-point line --- the best a team could do with points
alone). The first finding is negative and instructive: \textbf{text similarity alone carries almost no effort
signal}. The size-agnostic k-NN lands at MAE 3.83\,h --- statistically indistinguishable from the project-mean
baseline (3.82\,h) --- because text-similar neighbours are often \emph{size}-mismatched, so the estimate regresses to
the mean. Normalizing by size fixes it: the \textbf{size-aware} mode (draft points \texttimes{} the neighbours'
similarity-weighted hours-per-point rate) cuts MAE to \textbf{2.12\,h}, matching the story-points baseline
(2.10\,h). Unlike a regression line, it also returns the neighbour tasks as evidence the user can inspect, and
an honest uncertainty band that covers the true value 74\% of the time (against the \textasciitilde80\% its 10/90 band
nominally targets). In the delivered form the suggestion is size-agnostic (a draft being typed carries no
points yet); the size-aware mode activates whenever a pointed draft is estimated through the API.

\begin{longtable}{@{}lcccc@{}}
\caption{Estimation leave-one-out results (n = 43 completed tasks).}\label{tab:eval-estimation} \\
\tableheaderrow
\thcell{Predictor} & \thcell{MAE (h)} & \thcell{RMSE (h)} & \thcell{Within $\pm$25\%} & \thcell{Band coverage} \\
\midrule
\endfirsthead
\caption[]{Estimation leave-one-out results (continued).} \\
\tableheaderrow
\thcell{Predictor} & \thcell{MAE (h)} & \thcell{RMSE (h)} & \thcell{Within $\pm$25\%} & \thcell{Band coverage} \\
\midrule
\endhead
project-mean baseline & 3.82 & 5.22 & 42.9\% & --- \\
size-agnostic k-NN (text only) & 3.83 & 5.51 & 44.2\% & 0.488 (25/75) \\
story-points fit baseline & 2.10 & 2.58 & 62.8\% & --- \\
\textbf{size-aware k-NN (hours-per-point)} & \textbf{2.12} & \textbf{2.77} & \textbf{67.4\%} & \textbf{0.744 (10/90)} \\
\bottomrule
\end{longtable}

\section{Challenges \& design decisions}

Several non-obvious choices shaped the subsystem. Each has a clear reason:

\begin{itemize}
  \item \textbf{1536-dimensional embeddings, not the model's native 3072.} Gemini's embedding model supports Matryoshka
  truncation --- the first 1536 dimensions of a 3072-vector are themselves a usable embedding. Halving the
  dimension halves storage and speeds up the ANN scan for almost no quality loss on this corpus. Because reduced
  Gemini vectors are not pre-normalized and cosine ANN needs unit vectors, we L2-normalize them in application
  code before storing.
  \item \textbf{RRF with k = 60 to fuse the arms.} Reciprocal Rank Fusion~\cite{cormack2009rrf} combines the two ranked lists by rank position
  (\texttt{$\Sigma$ 1/(60 + rank)}) rather than by raw score, which avoids the problem that cosine similarity and
  \texttt{ts\_rank\_cd} are on different scales. The constant 60 is the standard damping value that keeps a single
  arm's top rank from dominating.
  \item \textbf{A cosine-only confidence gate at 0.5.} Refusing before generation is the cheapest way to avoid a
  hallucination, and 0.5 is the threshold below which retrieved content is too dissimilar to trust. This gate
  is effective but has a known limitation --- it ignores the lexical arm --- which we carry into the review (Section~4.10).
  \item \textbf{Why Gemini.} Using a managed provider~\cite{gemini} for both embeddings and generation avoided installing and
  operating an inference stack for a single sprint, and Gemini's \texttt{flash-lite} gave us a cheap, fast reranking
  model on its own quota bucket, separate from the main generation quota. The trade-off is a dependency on an
  external API and its rate limits --- visible in the evaluation note above.
  \item \textbf{A generated, stored \texttt{tsvector} instead of a maintained column.} Declaring the full-text column
  \texttt{GENERATED ALWAYS AS (...) STORED} means Postgres back-fills every existing row when the migration runs and
  keeps it current on every write, with no application code on the write path and no backfill script.
  \item \textbf{The outbox boundary, fire-and-forget.} Making index enqueues catch their own errors is a deliberate
  coupling decision: the AI subsystem depends on the domain (it reads tasks, epics, sprints) but the domain
  must not depend on it, so a failure in indexing can never appear as a failed task save.
\end{itemize}

\section{Sprint review}

The subsystem is the cleanest module in the codebase --- each concern (embedding, chunking, the outbox boundary,
retrieval, reranking, generation, access control, telemetry) is a single-responsibility class, the raw SQL is
isolated in two repositories, and the docstrings explain \emph{why} as well as \emph{what}. It delivers all seven
stories, and it is the only feature backed by measured results. The same review also documents its
verified gaps:

\begin{itemize}
  \item \textbf{Comment edits and deletes are not enqueued on the write path.} Only comment \emph{create} enqueues an index
  job; edits and deletes rely on the nightly reconcile. The consequence is the one property that undermines the
  ``never cite what does not exist'' promise: a \textbf{deleted comment stays retrievable and citable for up to \textasciitilde24 h}
  until reconcile runs, and the copilot can then cite content that is gone. This is the module's most
  important known issue and the first item to address --- copy the task producers and enqueue a
  \texttt{TASK\_COMMENT} upsert on edit and a delete on removal.
  \item \textbf{The confidence gate is cosine-only.} Because \texttt{sufficient} is derived solely from the top cosine score, a
  bare-keyword query that the \emph{lexical} arm answers perfectly but the \emph{dense} arm scores below 0.5 is still
  refused --- so the copilot's answer path does not fully inherit the hybrid keyword improvement demonstrated in Section~4.8. Making
  the gate consider ``any arm produced a hit above its own threshold'' would close this gap.
\end{itemize}

That these two gaps were found by the module's own review is the point: the same evaluation discipline that
measured the retrieval also showed where the answer path falls behind it. The remaining, smaller gaps from the same
review --- cancelled SSE streams that skip their telemetry receipt, the plaintext \texttt{CopilotQueryLog}, source-text
truncation shorter than the chunk, the reranker's scoring of omitted candidates, in-memory reindex/reconcile,
a missing fail-fast on the Gemini key, and the absence of unit tests around the pure functions (RRF fusion,
weighted percentiles, citation parsing) --- are collected with their fixes in Annex A. In priority order, the
two items above come first: enqueue comment edits and deletes so the copilot never cites content that is gone
(correctness), then make the confidence gate arm-aware so the answer path inherits the proven keyword improvement (UX).

\section{Sprint retrospective}

\textbf{Difficulties encountered.} The main challenge was working with raw SQL inside a
Prisma-managed codebase. Prisma cannot type pgvector or tsvector columns, so the embedding and
full-text search queries had to be written as parameterized raw templates and quarantined in
dedicated repositories. Ensuring these raw queries remained injection-safe while handling the
polymorphic \texttt{entityId} soft reference required careful attention. The asynchronous nature
of the indexing pipeline --- with its outbox boundary, retry backoff, and nightly reconcile --- also
made end-to-end verification more difficult than for synchronous modules.

\textbf{Lessons learned.} Building the offline evaluation harness early proved extremely useful: it
caught the keyword retrieval gap before the feature was delivered and validated the hybrid fix
quantitatively. The CRISP-DM framing also helped --- naming the phases made it explicit that
the evaluation harness is the quality gate playing the role of a trained model's validation curve,
which guided both design and review discussions.

\textbf{Improvements planned.} Enqueue comment edits and deletes on the write path so the
copilot never cites deleted content. Make the confidence gate arm-aware so the answer path
fully inherits the proven hybrid keyword improvement.

\section{Final cumulative class diagram}

Figure~\ref{fig:4-6} completes the picture: the whole system after six sprints, with the earlier sprints
compacted to their identifying anchor classes and the new \textbf{AI / RAG} package highlighted in yellow. The AI package
attaches to the rest in a deliberately loose manner. \texttt{DocumentEmbedding} and
\texttt{IndexOutbox} are attached to \texttt{Project} by their \texttt{projectId} scope key, and \texttt{CopilotQueryLog} to \texttt{User}; but the
link to the \emph{content} --- tasks, comments, epics, sprints, milestones --- is drawn as a dashed \textbf{soft reference},
because those \texttt{entityId} columns carry no foreign key (they are polymorphic across the five types). The only
inbound coupling is the write-path enqueue: the task, epic, sprint, and milestone services enqueue into
\texttt{IndexOutbox}, shown as dependencies, and that is the entire surface the AI subsystem exposes to the domain ---
fire-and-forget one-line calls that the domain can ignore. High cohesion inside, loose coupling outside.

% Figure 4.6 — render diagrams-src/fig_4_6.puml -> figures/fig_4_6.(png|pdf)
\reportsideways{figures/fig_4_6}{Final cumulative class diagram after six sprints.}{4-6}

\section{Sprint conclusion}

Sprint~6 delivered the platform's AI subsystem: a permission-scoped copilot with grounded,
citation-backed answers and an evidence-based task estimator, both built on a hybrid retrieval
engine whose quality is measured by an offline evaluation harness. The module integrates with
the domain through a deliberately loose outbox boundary and attaches to the cumulative data model
via soft references, keeping the AI subsystem self-contained while reading from all the project
content accumulated over the five preceding sprints.
